% Part:sets-functions-relations
% Chapter: size-of-sets
% Section: comparing-sizes

\documentclass[../../../include/open-logic-section]{subfiles}

\begin{document}

\olfileid{sfr}{siz}{car}

\olsection{د بېلو اندازو سټونه او د کانتور قضيه}

\begin{explain}
موږ دا مفکوره په دقيقه وينا کښې وړاندې کړې چې دوه سټونه يوه
اندازه لري۔ دا مفکوره هم په دقيقه وينا کښې وړاندې کولے شو چې يو
سټ له بل څخه کوچنے دے۔ زموږ د «له ... څخه کوچنے (يا همشمېره)
دے» تعريف به د سټونو تر منځ د !!a{bijection} پر ځاے، له لومړي
سټ څخه دويم ته يوه !!a{injection} وغواړي۔ که داسې تابع وي، د
لومړي سټ اندازه د دويم سټ له اندازې څخه کوچنۍ يا ورسره مساوي ده۔
په شهودي ډول، له يوه سټ څخه بل ته !!a{injection} دا ډاډ راکوي
چې د تابعې د قيمتونو سټ لږ تر لږه د تعريف د ساحې هومره
!!{element}s لري، ځکه د تعريف د ساحې هېڅ دوه !!{element}s د
قيمتونو د سټ يوه !!{element} ته نۀ نقشېږي۔
\end{explain}

\begin{defn}
$A$ له~$B$ څخه \emph{لوے نۀ دے}، چې $\cardle{A}{B}$ ئې ليکو،
هله او يوازې هله چې يوه !!a{injection} $f \colon A \to B$ وي۔
\end{defn}

څرګنده ده چې دا يوه انعکاسي او انتقالي اړيکه ده، خو متناظره نۀ
ده (دا د تمرين په توګه پرېښوول شوې ده)۔ يوه داسې مفکوره هم
معرفي کولے شو چې وايي يو سټ له بل څخه (په ټينګه) کوچنے دے۔

\begin{defn}
$A$ له~$B$ څخه \emph{کوچنے دے}، چې $\cardless{A}{B}$ ئې ليکو،
هله او يوازې هله چې يوه !!a{injection}~$f\colon A \to B$ وي،
خو هېڅ !!{bijection}~$g\colon A \to B$ نۀ وي؛ يعنې
$\cardle{A}{B}$ او $\cardneq{A}{B}$ وي۔
\end{defn}

څرګنده ده چې دا اړيکه ناانعکاسي او انتقالي ده۔ (دا د تمرين په
توګه پرېښوول شوې ده۔) په دې نښه ويلے شو چې يو سټ $A$ هله او
يوازې هله !!{enumerable} دے چې $\cardle{A}{\Nat}$ وي، او $A$
هله او يوازې هله !!{nonenumerable} دے چې
$\cardless{\Nat}{A}$ وي۔
\textbf{د ژباړې د نسخې سرچينيز شرط (PSSIZ-004):} وروستے
هم‌ارز بيان دا فرض کاروي چې هر نامتناهي سټ د طبيعي عددونو يوه
يو پر يو کاپي لري؛ دا د انتخاب له اصل پرته په عمومي سټ‌تيورۍ
کښې نۀ ثابتېږي۔ په دې سره دا بيا ويلے شو:
\oliflabeldef{sfr:siz:nen-alt:thm:nonenum-pownat}{%
\olref[sfr][siz][nen-alt]{thm:nonenum-pownat}
د دې مشاهدې په توګه چې
$\cardless{\Nat}{\Pow{\Nat}}$}{%
\olref[sfr][siz][nen]{thm:nonenum-pownat}
د دې مشاهدې په توګه چې $\cardless{\PosInt}{\Pow{\PosInt}}$}۔ په
حقيقت کښې، \citet{Cantor1892} ثابته کړه چې دا وروستۍ خبره
\emph{بشپړه عمومي} ده:

\begin{thm}[کانتور]\ollabel{thm:cantor}
د هر سټ $A$ دپاره $\cardless{A}{\Pow{A}}$۔
\end{thm}

\begin{proof}
نقشه $f(x) = \{x\}$ يوه !!a{injection}
$f \colon A \to \Pow{A}$ ده، ځکه که $x \neq y$ وي، نو د غړو له
مخې د برابرۍ په سبب $\{x\} \neq \{y\}$ هم وي، او نو
$f(x) \neq f(y)$۔ له دې امله $\cardle{A}{\Pow{A}}$ لرو۔

\begin{editorial}
که \olref[nen]{sec} موجوده وي، تفصيلي ثبوت وړاندې کوو؛ ګنې يو
چټک ثبوت وړاندې کوو چې د \olref[nen-alt]{sec} سره سمون خوري۔
\end{editorial}

\oliflabeldef{sfr:siz:nen:sec}{% 
  اوس به وښيو چې هېڅ !!a{surjective} تابع~$g\colon A \to
  \Pow{A}$ نۀ شي کېدے، نو يوه !!a{bijective} تابع خو لا پرېږده؛
  او له دې امله $\cardneq{A}{\Pow{A}}$۔ ځکه فرض کړئ چې
  $g\colon A \to \Pow{A}$۔ دا چې $g$ بشپړه ده، هر $x \in A$ د
  يوه فرعي سټ $g(x) \subseteq A$ سره نقشېږي۔ موږ ښودلے شو چې
  $g$ پر هدف سټ بشپړه نۀ شي کېدے۔ د دې دپاره يو فرعي سټ
  ~$\overline{A} \subseteq A$ تعريفوو چې د تعريف له مخې د~$g$
  په قيمتونو سټ کښې نۀ شي راتلے۔ کېږدئ:
  \[
  \overline{A} = \Setabs{x \in A}{x \notin g(x)}.
  \]
  ځکه چې $g(x)$ د هر $x \in A$ دپاره تعريف شوے، $\overline{A}$
  په څرګند ډول د~$A$ يو سم تعريف شوے فرعي سټ دے۔ خو دا د~$g$
  د قيمتونو په سټ کښې نۀ شي راتلے۔ يو خپل سری $x \in A$ ونيسئ؛
  موږ به وښيو چې $\overline{A} \neq g(x)$۔ که $x \in g(x)$ وي،
  نو $x \notin g(x)$ نۀ پوره کوي، او د~$\overline{A}$ د تعريف له
  مخې $x \notin \overline{A}$ لرو۔ که $x \in \overline{A}$ وي،
  د~$\overline{A}$ تعريفي خاصيت بايد پوره کړي؛ يعنې $x \in A$
  او $x \notin g(x)$ وي۔ دا چې $x$ خپل سری و، دا ښيي چې د هر
  $x \in A$ دپاره، $x \in g(x)$ هله او يوازې هله چې
  $x \notin \overline{A}$، او نو $g(x) \neq \overline{A}$۔
  \textbf{د ژباړې د نسخې سمون (OLSIZ-007):} انګرېزي ثبوت په
  دې عمومي پايله کښې ورودي د ټول سټ پر ځاے يوازې قطري فرعي سټ
  ته محدود کړے و؛ دلته ئې د مخکښې ټاکل شوي خپل سري ورودي ساحه
  راوګرځول شوه۔ په بله وينا، $\overline{A}$ د~$g$ په قيمتونو
  سټ کښې نۀ شي راتلے، او دا له دې فرض سره تناقض لري چې~$g$ پر
  هدف سټ بشپړه ده۔}{اوس يوازې دا ښودل پاتې دي چې
  $\cardneq{A}{\Pow{A}}$۔ د خلف دپاره فرض کړئ
  $\cardeq{A}{\Pow{A}}$؛ يعنې کومه !!{bijection}
  $g \colon A \to \Pow{A}$ شته۔ اوس دا په پام کښې ونيسئ:
  \[
    D = \Setabs{x \in A}{x \notin g(x)}
  \]
  په ياد ولرئ چې $D \subseteq A$، نو $D \in \Pow{A}$۔ ځکه چې
  $g$ يوه !!a{bijection} ده، کوم $y \in A$ شته چې $g(y) = D$۔
  خو اوس لرو:
  \[
    y \in g(y) \text{ هله او يوازې هله چې } y \in D \text{ هله او يوازې هله چې } y \notin g(y).
  \]
  دا يو تناقض دے؛ نو $\cardneq{A}{\Pow{A}}$۔}{}
\end{proof}

\begin{explain}
\oliflabeldef{sfr:siz:nen:thm:nonenum-pownat}{د
  \olref{thm:cantor} ثبوت د
  \olref[nen]{thm:nonenum-pownat} له ثبوت سره پرتله کول ګټور دي۔
  هلته مو وښودل چې د~$\PosInt$ د فرعي سټونو د هر لړ $Z_1$،
  $Z_2$، \dots دپاره د عددونو يو سټ~$\overline{Z}$ جوړولے شو
  چې په ډاډه توګه په لړ کښې نۀ وي۔ په لړ کښې نۀ و، ځکه چې د هر
  $n \in \PosInt$ دپاره، $n \in Z_n$ هله او يوازې هله چې
  $n \notin \overline{Z}$۔ په دې توګه تل کوم عدد د $Z_n$ او
  $\overline{Z}$ له يوه څخه !!a{element} وي خو د بل نۀ وي۔ دلته
  هماغه مفکوره تعقيبوو، خو شاخصونه~$n$ اوس د~$\PosInt$ پر ځاے
  د~$A$ !!{element}s دي۔ سټ $\overline{A}$ داسې تعريف شوے چې د
  هر $x \in A$ دپاره له~$g(x)$ څخه بېل وي، ځکه $x \in g(x)$
  هله او يوازې هله چې $x \notin \overline{A}$۔ بيا هم تل د~$A$
  يو !!a{element} شته چې له $g(x)$ او $\overline{A}$ څخه د يوه
  !!a{element} وي خو د بل نۀ وي۔ او لکه څنګه چې له دې امله
  $\overline{Z}$ په لړ $Z_1$، $Z_2$، \dots کښې نۀ شي راتلے،
  $\overline{A}$ د~$g$ په قيمتونو سټ کښې نۀ شي راتلے۔}{}

\oliflabeldef{sfr:siz:nen-alt:thm:nonenum-pownat}{د
  \olref{thm:cantor} ثبوت د
  \olref[nen-alt]{thm:nonenum-pownat} له ثبوت سره پرتله کول ګټور
  دي۔ هلته مو وښودل چې د~$\Nat$ د فرعي سټونو د هر لړ $N_0$،
  $N_1$، $N_2$، \dots دپاره د عددونو يو سټ~$D$ جوړولے شو چې په
  ډاډه توګه په لړ کښې نۀ وي۔ په لړ کښې نۀ و، ځکه چې د هر
  $n \in \Nat$ دپاره $n \in N_n$ هله او يوازې هله چې
  $n \notin D$۔ دلته هماغه مفکوره تعقيبوو، خو شاخصونه~$n$ اوس
  د~$\Nat$ پر ځاے د~$A$ !!{element}s دي۔ سټ $D$ داسې تعريف
  شوے چې د هر $x \in A$ دپاره له~$g(x)$ څخه بېل وي، ځکه
  $x \in g(x)$ هله او يوازې هله چې $x \notin D$۔}{}

دا ثبوت د رسل د پاراډوکس له ثبوت،
\olref[sfr][set][rus]{thm:russells-paradox}، سره هم د پرتله کولو
وړ دے۔ په حقيقت کښې، د کانتور قضيه د رسل د خپل پاراډوکس الهام و۔
\end{explain}

\begin{prob}
  وښيئ چې د هر سټ~$A$ دپاره هېڅ !!a{injection}
  $g\colon \Pow{A} \to A$ نۀ شي کېدے۔ لارښوونه: فرض کړئ
  $g\colon \Pow{A} \to A$ !!{injective} ده۔
  $D = \Setabs{g(B)}{B \subseteq A \text{ او } g(B) \notin B}$
  په پام کښې ونيسئ۔ کېږدئ $x = g(D)$۔ د تناقض د راايستلو دپاره
  دا حقيقت وکاروئ چې $g$ !!{injective} ده۔
\end{prob}

\end{document}
